\section{Positive level currents and an intrinsic energy disintegration}
\label{sec:r11-level-disintegration}

The construction in this section applies to every allowed original
background.  It introduces positive currents carried by the level sets
of a bounded difference of classical psh functions.  They need not be
closed.  Their exact differential defect is retained, and their
Lebesgue integral recovers the energy current.  In particular no
arbitrary-current pullback, auxiliary residue field, or analytic
PS-relation quotient occurs in the construction.

Keep the original finite extension $K/\mathbf Q_p$ and its absolute
value.  Let $X$ be good, Hausdorff, boundaryless, quasi-smooth and pure
of dimension $n$, and let $U\subset X$ be ordinary open.  Every test
below is an original compactly supported smooth CLD test, of its
complementary bidegree.  Let $Q$ be an arbitrary symmetric
middle-positive, separately $d'$- and $d''$-closed current of degree
$(p,p)$, where $0\le p<n$.  Write $D=d'd''$.  A bounded difference
means a function $f=a-b$, with $a,b$ locally bounded classical psh;
all assertions are local, and finitely many such presentations may
be used on a compact test support.

The inputs are the frozen bounded weighted BT calculus, finite
polarization, strict Borel maximum locality, and coefficient-Radon
control on all original tests.  These are
\texttt{INTERNAL\_INPUT}.  In particular $D(fQ)$ is a signed
coefficient-Radon current, defined by linearity from the two bounded
classical potentials.

\subsection{The bounded-difference chain rule}

\begin{lemma}[Smooth scalar calculus on bounded differences]
\label{lem:r11-bounded-difference-chain}
The distribution
\[
 \Gamma_Q(f,g)
 =\frac12\{D(fgQ)-fD(gQ)-gD(fQ)\}
 \tag{LC1}
\]
is a symmetric coefficient-Radon current of degree $(p+1,p+1)$.
It is bilinear, and $\Gamma_Q(f):=\Gamma_Q(f,f)$ is
middle-positive.  On every relatively compact witness stage,
for a smooth real function $F$ on an interval containing the range
of $f$, one has
\[
 D(F(f)Q)=F'(f)D(fQ)+F''(f)\Gamma_Q(f),
 \qquad
 \Gamma_Q(F(f),g)=F'(f)\Gamma_Q(f,g).
 \tag{LC2}
\]
All multipliers use the given actual Borel values.  The formulas
are independent of the psh presentation of $f,g$ and of every
actual moment refinement.
\end{lemma}

\noindent\textbf{Status: \texttt{DIRECT\_COROLLARY}} of the
frozen weighted bounded calculus.  The passage from top-degree
energy measures to lower-degree coefficient currents is included
in the following argument.

\begin{proof}
Take simultaneous smooth decreasing classical rows for the finitely
many psh summands on a compact regularization stage.  The smooth
identity for (LC1) is
\[
 \Gamma_Q(f,g)
 =\tfrac12(d'f\wedge d''g+d'g\wedge d''f)\wedge Q.
\]
For $f=g$ the first factor is strongly positive, so the current is
middle-positive.  The two weighted terms in (LC1) converge by the
frozen weighted theorem, and the first term converges by bounded
convergence against the fixed coefficient measures of $Q$.
Thus (LC1) is the distributional limit of positive currents when
$f=g$.  The original positive-current coefficient theorem makes
it coefficient-Radon; polarization treats $f,g$.  This proof works
in every permitted bidegree, not only in top degree.

For completeness, a smooth function $H$ of finitely many bounded
psh variables is locally a difference of two smooth convex,
coordinatewise nondecreasing functions of those variables.  On a
slightly larger compact real box, add a sufficiently large positive
multiple of $\sum x_i^2$, followed by a sufficiently large linear
multiple of $\sum x_i$.  The resulting function and the added
function are convex and coordinatewise nondecreasing on that box.
Their compositions with the psh variables are bounded classical
psh functions.  Applied to decreasing rows, each such composition
decreases to the actual composition.  This gives, in particular,
bounded-difference presentations for $F(f)$ and all smooth scalar
coefficients needed in the smooth chain identities.

One can therefore pass every weighted term of those identities by
the frozen joint decreasing theorem after the indicated finite
decompositions.  For an energy term multiplied by a smooth scalar
coefficient, expand (LC1): the $D(f^2Q)$ term is a difference of
bounded one-potential BT currents by the same convex decomposition
of $f^2$, and its coefficient, as well as the coefficients of
$fD(fQ)$, has such a decomposition.  Every passage is consequently
a finite combination of the already proved weighted BT passages.
This proves both formulas in (LC2).  They involve only intrinsic
distributional operations and actual coefficient multiplication,
which proves independence of the chosen witnesses and refinements.
\end{proof}

\begin{definition}[The level curvature current]
\label{def:r11-level-current}
For a bounded difference $f$ and $t\in\mathbf R$ put
\[
 \mathsf L^Q_f(t)
 :=D(\max\{f,t\}\,Q)
                 -\mathbf1_{\{f>t\}}D(fQ).
 \tag{LC3}
\]
The first current is well defined: if $f=a-b$, then
$\max\{f,t\}=\max\{a,b+t\}-b$ is again a bounded
difference.  Formula (LC3) is the definition, not an assertion
that a Borel-restricted closed current is closed.
\end{definition}

\begin{theorem}[Positivity, carriers, and the differential defect]
\label{thm:r11-positive-level-current}
For every real $t$, $\mathsf L^Q_f(t)$ is a middle-positive
coefficient-Radon current of degree $(p+1,p+1)$ and is carried by
the actual Borel set $\{f=t\}$.  More precisely,
\[
 \mathbf1_{\{f\ne t\}}\mathsf L^Q_f(t)=0.
 \tag{LC4}
\]
Its separate derivatives are exactly
\[
 \begin{split}
 d'\mathsf L^Q_f(t)
       &=-d'\bigl(\mathbf1_{\{f>t\}}D(fQ)\bigr),\\
 d''\mathsf L^Q_f(t)
       &=-d''\bigl(\mathbf1_{\{f>t\}}D(fQ)\bigr).
 \end{split}
 \tag{LC5}
\]
No vanishing of these derivatives is asserted when $p+1<n$.
The construction is local in the original ordinary topology and
commutes with every actual refinement of the same functions and tests.
\end{theorem}

\noindent\textbf{Status: \texttt{NEW\_PROVED}}.

\begin{proof}
Choose a smooth convex $\psi_\epsilon:\mathbf R\to\mathbf R$
which is zero on $(-\infty,0]$, has derivative in $[0,1]$, and
has derivative one on $[\epsilon,\infty)$.  It can be obtained by
integrating twice a nonnegative smooth function of integral one
supported in $(0,\epsilon)$.  It satisfies
\[
 |\psi_\epsilon(s)-s_+|\le\epsilon,
 \qquad
 \psi_\epsilon'(s)\longrightarrow\mathbf1_{\{s>0\}}
                   \quad\text{at every real }s.
\]
The chain identity (LC2) gives
\[
 D(\psi_\epsilon(f-t)Q)
        -\psi_\epsilon'(f-t)D(fQ)
            =\psi_\epsilon''(f-t)\Gamma_Q(f)\ge0.
 \tag{LC6}
\]
On a compact original test support the first term tends to
$D((f-t)_+Q)$ by uniform convergence and the fixed coefficient
Radon bounds for $Q$.  The second tends to
$\mathbf1_{\{f>t\}}D(fQ)$ by bounded convergence against the
fixed signed coefficient measures of $D(fQ)$.  As $D(tQ)=0$,
the limiting distribution is (LC3).  Positivity is closed under
this convergence on all original tests, proving the claim.

To identify its carrier, use $\max\{f,t\}=\max\{a,b+t\}-b$.
Strict Borel maximum locality gives
\[
 \begin{split}
 \mathbf1_{\{f>t\}}D(\max\{f,t\}Q)
        &=\mathbf1_{\{f>t\}}D(fQ),\\
 \mathbf1_{\{f<t\}}D(\max\{f,t\}Q)&=0.
 \end{split}
\]
This is a coefficient-measure identity; neither strict Borel set
is being declared ordinary open.  Subtracting the second term
of (LC3) proves (LC4).  Applying $d'$ and $d''$ to (LC3),
and using $d'D=d''D=0$, proves (LC5).  Finally every term in
the definition already has the original all-test and refinement
meaning, so no additional descent is needed.
\end{proof}

\subsection{Coarea on all original tests}

\begin{theorem}[Lebesgue disintegration of the energy current]
\label{thm:r11-level-coarea}
For every bounded Borel function $\vartheta:\mathbf R\to\mathbf R$
and every original compact smooth complementary test $\eta$,
\[
 \int_{\mathbf R}\vartheta(t)
                 \langle\mathsf L^Q_f(t),\eta\rangle\,dt
       =\langle\vartheta(f)\Gamma_Q(f),\eta\rangle.
 \tag{LC7}
\]
The integration is local on the test support: levels outside a
compact interval containing its nearby $f$-range give zero.
In particular,
\[
 \int_{\mathbf R}\mathsf L^Q_f(t)\,dt=\Gamma_Q(f),
 \qquad
 \mathbf1_{\{f=t\}}\Gamma_Q(f)=0\quad(t\in\mathbf R).
 \tag{LC8}
\]
The second identity holds for every level, including levels
charged by $D(fQ)$.

For a fixed positive trace on a compact moment stage, the resulting
positive scalar measures satisfy the full kernel formula
\[
 \int_B \mathsf L^Q_f(t)(E)\,dt
       =\Gamma_Q(f)(E\cap\{f\in B\})
 \tag{LC9}
\]
for Borel $B\subset\mathbf R$ and Borel $E$ in that stage.
Here the notation in (LC9) refers to the chosen scalar trace;
finite polarization recovers every original coefficient pairing.
\end{theorem}

\noindent\textbf{Status: \texttt{NEW\_PROVED}}.

\begin{proof}
First take $\vartheta\in C_c^\infty(\mathbf R)$ and set
\[
 F(s)=\int_{\mathbf R}\vartheta(t)(s-t)_+\,dt.
\]
Then $F''=\vartheta$ and
$F'(s)=\int_{t<s}\vartheta(t)\,dt$.  The function $F$ is
smooth.  On a compact test support the integrands in the
potential term of (LC3) are uniformly bounded, while the
restricted-curvature term is controlled by the fixed total
variations of $D(fQ)$.  Fubini and distributional differentiation
therefore give
\[
 \int\vartheta(t)\mathsf L^Q_f(t)\,dt
      =D(F(f)Q)-F'(f)D(fQ)
      =\vartheta(f)\Gamma_Q(f).
\]
The last equality is (LC2).

For a positive compact smooth trace test, the map
$t\mapsto\langle\mathsf L^Q_f(t),\eta\rangle$ is
nonnegative and Borel: the maximum term is continuous in $t$
on each compact support, and the restricted signed-measure term
is Borel by ordinary dominated integration.  Its integral is
finite on every relevant compact interval by the proved formula
with a smooth cutoff equal to one there.  Uniqueness of finite
Radon measures on $\mathbf R$ extends the formula from smooth
$\vartheta$ to bounded Borel $\vartheta$.  A finite collection
of positive trace tests dominates all coefficient variations
on a compact stage, so finite polarization gives (LC7) for
every original $\eta$.

For a fixed scalar trace, weak Borel measurability of the positive
Radon measures $\mathsf L^Q_f(t)$ implies measurability of their
values on Borel sets.  This can also be checked directly by
approximating open-set indicators from a countable continuous
core on a relatively compact affinoid stage and then using the
monotone-class theorem.  Formula (LC7), first for compact
continuous scalar weights on the stage and then by Radon
uniqueness, gives (LC9).  These arguments are local on finitely
many such stages and do not impose global second countability
on $X$.

Taking $B=\{t\}$ in (LC9) proves the zero-level energy assertion,
because Lebesgue measure has no atoms.  Taking $\vartheta=1$
on the local range proves the first identity in (LC8).
\end{proof}

\begin{corollary}[A localized bounded-variation energy density]
\label{cor:r11-energy-bv-density}
Suppose now that $p=n-1$, so the level currents are measures.
For an original smooth compact $\chi\ge0$ define the finite
signed measure on $\mathbf R$
\[
 \nu_{\chi,f,Q}
  :=f_*\bigl(\chi D(fQ)+\Gamma_Q(\chi,f)\bigr).
 \tag{LC10}
\]
Then $\nu_{\chi,f,Q}(\mathbf R)=0$, and for every $t\in\mathbf R$
\[
 k_{\chi,f,Q}(t):=\int\chi\,d\mathsf L^Q_f(t)
       =\nu_{\chi,f,Q}(( -\infty,t])\ge0.
 \tag{LC11}
\]
Thus $k$ is a right-continuous compactly supported BV function,
$\operatorname{Var}(k)\le|\nu_{\chi,f,Q}|(\mathbf R)$, and
\[
 f_*(\chi\Gamma_Q(f))=k_{\chi,f,Q}(t)\,dt,
 \qquad
 \|k_{\chi,f,Q}\|_\infty
                     \le|\nu_{\chi,f,Q}|(\mathbf R).
 \tag{LC12}
\]
\end{corollary}

\noindent\textbf{Status: \texttt{NEW\_PROVED}}.

\begin{proof}
The compact Stokes identity gives
$\int\chi D(fQ)=-\int\Gamma_Q(\chi,f)$, hence total mass
zero in (LC10).  Apply the first-derivative chain identity to the
one-sided smooth truncations used in (LC6).  Their uniform
convergence and the pointwise convergence of their derivatives
give
\[
 \int\chi D(\max\{f,t\}Q)
       =-\int_{\{f>t\}}\Gamma_Q(\chi,f).
\]
All limits here concern the fixed finite mixed energy measure;
its total variation is finite by Cauchy--Schwarz.  Subtracting
the strict curvature term of (LC3) yields
$k(t)=-\nu_{\chi,f,Q}((t,\infty))$, which is (LC11).
The support of $\chi$ and its differential is compact, and
$f$ is bounded on a neighborhood of that support, so $\nu$
and $k$ have the claimed compact support.  The BV and norm
bounds are the elementary cumulative-measure bounds.  Finally
(LC7) identifies the pushforward density with $k$.
\end{proof}

\begin{proposition}[A nonclosed level current on the original two-torus]
\label{prop:r11-nonclosed-level-example}
The conclusion that every positive level current in
Theorem~\ref{thm:r11-positive-level-current} is separately closed is
\texttt{FALSE}, even when the potential is smooth and classical and
the background is the original integration current.  This failure
occurs in every dimension $n\geq2$.
\end{proposition}

\noindent\textbf{Status: \texttt{NEW\_PROVED}} counterexample.

\begin{proof}
First take $X=(\mathbf G_{m,K}^2)^{\mathrm{an}}$, set
$t_i=-\log|T_i|$, and work on the ordinary domain
$U=\{|t_1|<1,\ |t_2|<1\}$.  Put
\[
 Q=[U],\qquad f=t_1^2+t_2,\qquad
 B_1=d't_1\wedge d''t_1.
 \tag{LC13}
\]
The function $f$ is bounded, original smooth and classical psh:
its convex moment expression has positive semidefinite Hessian.
One has $D(fQ)=2B_1\wedge[U]$.  The distributional derivative of
the half-space indicator in the two original moment coordinates
therefore gives
\[
 d'\mathsf L_f^Q(0)
 =-2\delta(t_1^2+t_2)\,d't_2\wedge B_1\wedge[U].
 \tag{LC14}
\]
This can be verified directly by the smooth one-sided functions
from (LC6): their derivatives converge to the strict indicator,
and ordinary real integration by parts in the moment variables
gives the displayed hypersurface distribution.

Choose $\rho\in C_c^\infty((-1,1)^2)$ nonnegative and positive
on a nonempty segment of $\{(s,-s^2):|s|<1/2\}$.
The original smooth form $\rho(t_1,t_2)d''t_2$ has compact
support in $U$: the inverse image of the compact moment support
is a compact polyannulus.  Original tropical calibration of the
two-torus has multiplicity one, so its pairing with (LC14),
up to the fixed orientation convention, has absolute value
\[
 2\int_{-1}^{1}\rho(s,-s^2)\,ds>0.
 \tag{LC15}
\]
Thus $d'\mathsf L_f^Q(0)\ne0$.  By symmetry the double-prime
defect is nonzero as well.  In higher dimensions take the product
with $n-2$ original open annuli, keep $f$ independent of the extra
coordinates, and wedge the test with a positive compact moment
volume form in those coordinates.  Its nonzero pairing is the
product of (LC15) with a strictly positive real integral.  This
uses actual original tests in every dimension; no delta-test
enlargement is used.
\end{proof}

\paragraph{Closedness and the relation to the general target.}
The currents $\mathsf L^Q_f(t)$ provide an actual level
disintegration of energy, rather than a disintegration of an
arbitrary BT curvature measure into closed backgrounds.
Equation (LC5) is indispensable below top degree.  In particular
one cannot insert each level current into the frozen closed-background
calculus and integrate without separately resolving that derivative
defect.  Neither general OPEN-GQ nor general Borel regularity of
the raw relative BT capacity is inferred from this construction.
